% 392_Darstellungen_Trennung_En_ch.tex
% Johann Pascher, 11 Oct. 2026
% Where the representations of gravity part ways

\providecommand{\BM}{\textbf{[B]}}
\providecommand{\KM}{\textbf{[K]}}
\providecommand{\SM}{\textbf{[S]}}
\providecommand{\XM}{\textbf{[X]}}
\providecommand{\QM}{\textbf{[Q]}}

\chapter*{Where the Representations of Gravitation Part Ways}
\addcontentsline{toc}{chapter}{Where the Representations of Gravitation Part Ways}

\begin{abstract}
Doc.~389 treats space curvature, change of mass or clock rate, and fractal path lengthening as three
representations of gravity that must be equivalent by virtue of the duality $\tilde T\cdot m=1$. This
document examines how far that equivalence reaches. The temporal half $g_{00}=(m/m_0)^2$ is contained
in every admissible representation and alone carries free fall, redshift and clock rates. Pure mass
variation, in which clock rates and lengths both scale with $1/m$, is a conformal rescaling of flat
space and has a vanishing Weyl tensor. The flat Friedmann universe is conformally flat as well; there
the mass reading is exact and passes light curves, background temperature and the Tolman test,
whereas path lengthening acting on the photon alone fails all three. The field of a star, by
contrast, carries Weyl curvature, whose component $C_{trtr}=-2GM/r^3$ is the tidal term. There the
flat representation needs the spatial half with the opposite sign for rulers, a co-moving vector
field for rotating masses, and a tensor part for gravitational waves. The speed-of-light limit follows from the torus spectrum: massless modes travel at $c$, massive ones
are always slower; a conformal rescaling changes neither $c$ nor $\varepsilon_0$ and $\mu_0$ locally.
The mass reading is thus the
complete representation of the conformally flat part; the Weyl part is what has to be supplied in
addition. The Alcock--Paczyński test does not separate expansion and the mass reading, but with the
DESI data it excludes the pure exponential form $1+z=e^{\xi x}$ of the distance law.
\end{abstract}

% ============================================================
\section*{Status markers}
% ============================================================
\begin{center}
\begin{tabular}{@{}lp{11cm}@{}}
\textbf{[B]} & Algebraically proved (symbolic computation in the verification script). \\[2pt]
\textbf{[K]} & Numerically confirmed (verification script, comparison with measurements). \\[2pt]
\textbf{[S]} & Setting, open or speculative. \\[2pt]
\textbf{[X]} & Excluded by measurement. \\[2pt]
\textbf{[Q]} & Taken from the literature, with source. \\
\end{tabular}
\end{center}
Verification scripts: \texttt{2/python/Dok392\_Skripte/} (eight scripts, 81/81). Measured values and
sources are given in the script headers and in the bibliography.

% ============================================================
\section*{Symbols}
% ============================================================
\begin{center}
\small
\begin{tabular}{@{}p{3.3cm}p{11.2cm}@{}}
\toprule
$\xi=4/30000$ & the single parameter of FFGFT (geometry T$^4$/Z$_3$) \\
$\tilde T$, $m$, $m_0$ & time field, position-dependent mass, reference mass; duality $\tilde T\cdot m=1$ \\
$c$ & speed of light, $c=1$ in the basic form \\
$G$, $M$, $r$, $\vec x$ & gravitational constant, mass of the central body, distance, position vector \\
$\Phi$, $U$ & Newtonian potential; $U=GM/(rc^2)=-\Phi/c^2$, dimensionless \\
$g_{\mu\nu}$ & metric; $g_{00}$ temporal component (clock rate), $g_{0i}$ mixed (dragging),
$g_{ij}$ spatial components ($i,j=1,2,3$) \\
$\eta_{\mu\nu}$, $\Omega$ & flat Minkowski metric; conformal factor, $\Omega=m/m_0$ for pure mass
variation \\
$\gamma$, $\beta$ & PPN parameters: space curvature per unit mass; non-linearity of $g_{00}$ \\
$\eta=4\beta-\gamma-3$ & Nordtvedt parameter (not to be confused with $\eta_{\mu\nu}$ or the conformal
time $\eta$) \\
$m_G/m_I$ & ratio of gravitational to inertial mass (lunar laser ranging) \\
$C_{trtr}$ & component of the Weyl tensor, the tidal part of curvature \\
$n$ & refractive index of path lengthening, $n=1+2U$ \\
$\vec J$, $I$, $R$ & angular momentum, moment of inertia and radius of the rotating body \\
$\vec u$ & velocity field with which path lengthening is dragged along \\
$\vec w$ & magnetoelectric coupling in the medium form, $\vec w=-g_{0i}/g_{00}$ \\
$\dot P$ & rate of change of the orbital period of a binary \\
$z$, $a$, $a_0$ & redshift; scale factor, or $m/m_0$ in the mass reading; its present value \\
$\eta$ (cosmic), $H$, $H_0$ & conformal time; Hubble rate $H=a'/a^2$; its present value \\
$b$, $\beta$ (in $T(z)$), $k$ & exponents of light-curve stretching $(1+z)^b$, of the background
temperature $T\propto(1+z)^{1-\beta}$ (not the PPN parameter) and of surface brightness
$(1+z)^{-k}$ \\
W, T, M & cosmic readings: photon only; clocks in addition; clocks and rulers in addition \\
$d_L$, $d_A$ & luminosity and angular-diameter distance \\
$D_M$, $D_H$, $r_d$ & transverse comoving distance, Hubble distance $D_H=c/H$, sound-horizon length
(BAO ruler) \\
$x$ & path length in the law $1+z=e^{\xi x}$ (Docs.~280, 308) \\
$E$, $p$, $v$ & energy, momentum, group velocity $v=dE/dp$ of a mode \\
$x_4$, $R_4$, $k_4$, $\lambda_4$ & coordinate and radius of the compact fourth direction, winding number
($m=k_4/R_4$), wavelength with $\lambda_4\cdot m=2\pi$ \\
$\varepsilon$, $\mu$, $\varepsilon_0$, $\mu_0$ & effective permittivity and permeability in the field;
electric and magnetic constants \\
$A$, $B$ & metric functions in the static field, $g_{00}=-A$, $g_{ij}=B\,\delta_{ij}$ \\
$\chi^2$ & sum of squared deviations, measured in standard deviations \\
\bottomrule
\end{tabular}
\end{center}

% ============================================================
\section{The question}
% ============================================================

Doc.~389 reads gravity in three representations: as space curvature as in general relativity, as
change of mass or clock rate in flat space (Doc.~078), and as fractal path lengthening (Docs.~026,
041, 182, 308). It already shows there that the change of mass alone carries only the temporal half
and yields the light deflection of general relativity only together with path lengthening. It
remained open whether the three representations are equivalent beyond the static weak field, and if
not, where exactly they part ways. This document computes both sides: what excludes a
representation, and what speaks for the mass reading.

% ============================================================
\section{The temporal half as the common core}
% ============================================================

In the mass representation $g_{00}=(m/m_0)^2$ with $m=m_0\,e^{\Phi/c^2}$ (Doc.~389). Free fall follows
exactly from the mass gradient, $\ddot{\vec x}=-c^2\nabla\ln m=-\nabla\Phi$ \BM, and since
$\nabla\ln m$ is the same for every mode, the weak equivalence principle is built in. The same function
determines the redshift and the clock rate. The measurements of this half are:
\begin{center}
\begin{tabular}{@{}lll@{}}
\toprule
Measurement & Expectation from $g_{00}$ & measured \\
\midrule
Pound--Rebka, 22.5\,m & $2.45\times10^{-15}$ & $(2.57\pm0.26)\times10^{-15}$ \QM \\
optical clocks, 33\,cm & $3.6\times10^{-17}$ & $(4.1\pm1.6)\times10^{-17}$ \QM \\
GPS, clock-rate part per day & $+45.7\,\mu$s & corrected in operation \\
Gravity Probe A, Galileo 5/6 & clock-rate part & confirmed to $7\times10^{-5}$ and $2.5\times10^{-5}$ \QM \\
\bottomrule
\end{tabular}
\end{center}
\KM. This half is not an alternative to space curvature; it is the same in every admissible
representation, and in general relativity too $g_{00}$ is the clock rate. The mass reading describes
it literally: in a static field the frequency of the photon stays constant along the whole path, and
what differs are the atoms at the two places (Doc.~312). With $\tilde T\cdot m=1$ this is a single
quantity.

% ============================================================
\section{What the weak field excludes}
% ============================================================

In the parametrised weak field the representations differ in $\gamma$ (spatial half) and $\beta$
(second order of $g_{00}$). In the mass representation $\beta$ follows from the form of $m(\Phi)$: the
exponential form gives $g_{00}=1-2U+2U^2$, hence $\beta=1$, the linear form $g_{00}=1-2U+U^2$, hence
$\beta=\tfrac12$ \BM. Four independent measurements separate them:
\begin{center}
\small
\begin{tabular}{@{}lrrcccc@{}}
\toprule
Representation & $\gamma$ & $\beta$ & Light & Mercury & GP-B geodetic & Lunar laser \\
\midrule
Space curvature (GR) & 1 & 1 & $1.751''$ & $42.98''$ & 6604 & $\eta=0$ \\
Clock rate (exp.) + path lengthening & 1 & 1 & $1.751''$ & $42.98''$ & 6604 & $\eta=0$ \\
Clock rate (linear) + path lengthening & 1 & $\tfrac12$ & $1.751''$ & $50.1''$ \XM & 6604 & $\eta=-2$ \XM \\
clock rate only & 0 & 1 & $0.876''$ \XM & $14.3''$ \XM & 2201 \XM & $\eta=1$ \XM \\
conformally flat (Doc.~143) & $-1$ & $\tfrac12$ & $0$ \XM & $-7.2''$ \XM & $-2201$ \XM & $\eta=0$ \\
\bottomrule
\end{tabular}
\end{center}
\KM. Mercury in arcseconds per century, geodetic precession in mas per year, $\eta=4\beta-\gamma-3$.
The measured values are $\gamma-1=(2.1\pm2.3)\times10^{-5}$ (Cassini), the geodetic precession
$6601.8\pm18.3$\,mas/year (Gravity Probe B), and for lunar laser ranging
$(m_G/m_I)_E-(m_G/m_I)_M=(-0.8\pm1.3)\times10^{-13}$ \QM; with the difference of the self-energies of
Earth and Moon this gives $\eta=(1.8\pm2.9)\times10^{-4}$ without correction for composition \KM.
The change of mass alone thus fails also in a test without light, the geodetic precession, by about
240 standard deviations. The linear form of the mass change fails independently at Mercury and in
lunar laser ranging; the exponential form, carried as a setting in Doc.~389, is thereby fixed from
two sides.

Mixed forms are admissible. If curvature carries the share $w$ and path lengthening the share $s$ of
the spatial half, then $\gamma=w+s$, and the condition reads $w+s=1$ \KM. The freedom of
representation therefore lies in how the spatial half is divided, not in its size.

% ============================================================
\section{Conformally flat or Weyl-curved}
% ============================================================

If clock rates and lengths both scale with $1/m$, the metric is a conformal rescaling of flat space,
$g_{\mu\nu}=\Omega^2\eta_{\mu\nu}$ with $\Omega=m/m_0$. For every such metric the Weyl tensor vanishes
\BM. Pure mass variation therefore carries exactly the conformally flat geometry, and light travels in
it as in flat space.

\textbf{The universe at large.} The flat Friedmann universe can be written in conformal time as
$a(\eta)^2\eta_{\mu\nu}$; its Weyl tensor vanishes as well \BM. In the static frame the masses run as
$m(\eta)=m_0\,a(\eta)$, and $1+z=m_{\text{today}}/m_{\text{emission}}=a_0/a$
(Wetterich~\cite{Wetterich2013}) \QM. Because the cosmic geometry is conformally flat, the mass reading
is there not merely observationally equivalent but exactly the same geometry in other variables.

\textbf{The field of a star.} For the isotropic weak field with $\gamma=1$ the computation gives to
first order
\begin{equation}
C_{trtr}=-\frac{2GM}{r^3}\neq0\quad\BM,
\end{equation}
exactly the tidal term. In vacuum outside a mass the curvature is pure Weyl curvature. A conformal
rescaling cannot produce it.

\textbf{The sign of rulers.} Deep in the potential $m/m_0=1-U$. A length that scales with $1/m$, like
the Compton or Bohr length, appears longer by $1+U$ from the distant reference. The spatial half with
$\gamma=1$, by contrast, requires proper rulers that appear shorter by $1-U$ \BM. The signs are
opposite; this is the counter-running behaviour that Doc.~389 records for $\lambda_4\cdot m=2\pi$. In
the field of a star the spatial half therefore does not follow from the mass scaling, whereas in the
Friedmann universe it does.

% ============================================================
\section{Cosmos: which static reading is equivalent}
% ============================================================

For redshift without expansion there are three static readings (Docs.~280, 312, 391): path
lengthening or energy loss acting on the photon alone (W), the same with clocks booked along through
$\tilde T\cdot m=1$ (T), and the full mass reading, in which clocks and rulers are booked along (M).
Doc.~312 names four discriminators whose forward demonstration was open. Three of them can be
computed directly:
\begin{center}
\begin{tabular}{@{}lccc@{}}
\toprule
& Light curves $(1+z)^b$ & $T(z)\propto(1+z)^{1-\beta}$ & Tolman $(1+z)^{-k}$ \\
\midrule
Measurement \QM & $b=1.003\pm0.011$ & $\beta=-0.007\pm0.027$ & $k=2.59\pm0.17$ (R), $3.37\pm0.13$ (I) \\
W: photon only & $b=0$ \XM & $\beta=1$ \XM & $k=1$ \XM \\
T: + clocks & $b=1$ & $\beta=0$ & $k=2$ \XM \\
M: + clocks + rulers & $b=1$ & $\beta=0$ & $k=4$ \\
\bottomrule
\end{tabular}
\end{center}
\KM. The light curves exclude reading W by about 90 standard deviations, the background temperature
by about 37. In the Tolman test the observed values lie above $k=2$, although the evolution of the
stellar populations makes early galaxies brighter and raises the true exponent further; only $k=4$
is compatible. Equivalent to expansion is therefore exactly reading M. In it the rulers are booked
along by themselves, because atomic lengths scale with $1/m$ ($\lambda_4\cdot m=2\pi$, Doc.~314). Path
lengthening acting on the photon alone (Doc.~280) is not sufficient as a cosmic reading; it holds as a
description of the same redshift within the mass reading, not as an independent alternative.

The fourth discriminator, the Alcock--Paczyński test, compares the radial with the transverse extent of
an object. Computing in flat space with $m(t)=m_0a(t)$, with lengths and periods proportional to $1/m$
and conserved photon number, gives the angular-diameter and luminosity distances and the AP ratio
exactly as in the FRW model with the same $a$; in particular $d_L=(1+z)^2d_A$ \BM. Surface brightness
falls as $(d_A/d_L)^2$; the Tolman exponent $k=4$ is therefore the same statement as this reciprocity,
and readings W and T violate it with $d_L/d_A=(1+z)^{1/2}$ and $(1+z)$ respectively. The AP test
therefore does not separate expansion and the mass reading; it tests the course of $H(z)$, that is,
the distance law.

For the law $1+z=e^{\xi x}$ of Docs.~280 and 308 the test can be carried out if $x$ is read as the
coordinate distance of reading M. Then $a(\eta)=e^{H_0(\eta-\eta_0)}$, hence $H(z)=H_0(1+z)$, and the
ratio of transverse to radial distance becomes
\begin{equation}
\frac{D_M}{D_H}=(1+z)\ln(1+z).
\end{equation}
It contains neither $H_0$ nor the sound-horizon length $r_d$; the caveat P20, which Doc.~267 states for
the BAO distances, therefore does not apply. DESI DR2~\cite{DESI2025} measures this ratio in six
redshift ranges \QM. Up to $z\approx1.5$ the law lies within three standard deviations, and at
$z=0.51$ it matches the measured value exactly. For the Lyman-$\alpha$ forest at $z=2.33$ it gives
$4.01$ against a measured $4.52\pm0.10$, that is $5.3$ standard deviations too low; over all six values
$\chi^2=37.3$, whereas flat $\Lambda$CDM with $\Omega_m=0.2975$ gives $6.0$ \KM. In this reading the
pure exponential form is therefore excluded \XM. The degeneracy that Doc.~267 assumes for the BAO
angles requires a relation $z(d)$ that departs from the exponential law at high redshift; no such
relation has been derived in FFGFT \SM.

Equivalence alone does not decide between expansion and the mass reading. In favour of the mass
reading speaks the simpler structure of settings: the cosmological constant appears in it as an
artefact of the reading and is not set (Docs.~308, 312) \SM.

% ============================================================
\section{Rotating masses}
% ============================================================

Frame dragging by a rotating mass resides in $g_{0i}$ (Doc.~312). A static pair of scalar fields, mass
$m(x)$ and refractive index $n(x)$, has $g_{0i}=0$: the mass density of a rigidly rotating sphere is
stationary like that of a sphere at rest, the angular momentum $J$ does not appear in a scalar, and the
gravitomagnetic field $\nabla\times(\vec J\times\vec x/r^3)$ is a dipole field that only a vector field
can carry \BM. Gravity Probe B measures frame dragging as $37.2\pm7.2$\,mas/year against a prediction
of about 39, and LAGEOS confirms it to about 10\,\% \QM; absent frame dragging is thus excluded by 5
and about 10 standard deviations respectively \KM.

The flat representation reaches frame dragging if path lengthening is carried along like a moving
medium. The Gordon metric of a medium with index $n$ and velocity $\vec u$ has
$\tilde g_{0i}=-(1-n^{-2})\,u_i/c\approx-4U\,u_i/c$ for $n=1+2U$; equating this with the GR value
$g_{0i}=-2G(\vec J\times\vec x)_i/(c^3r^3)$ gives
\begin{equation}
\vec u=\frac{\vec J\times\vec x}{2Mr^2}\quad\BM.
\end{equation}
At the Earth's surface this is about $77$\,m/s at the equator, one sixth of the rotation velocity
($I/(2MR^2)$); for the Sun $70$\,m/s or $3.5\,\%$; the field falls off as $1/r$ \KM. This field is a
vector and does not follow from $m(x)$. In the curvature representation the dragging is contained,
because FFGFT adopts the field equations (Doc.~312); whether the structure of the carrier itself
supplies the co-moving field of the flat representation is open \SM.

% ============================================================
\section{Gravitational waves}
% ============================================================

Gravitational waves impose three requirements on every representation. First, they travel at the speed
of light: for GW170817 the gamma-ray signal arrived $1.74$\,s after the waves over about $40$\,Mpc,
$|c_{\text{gw}}/c-1|\lesssim10^{-15}$ \KM. Second, the spatial half acts equally on light and waves: the
Shapiro delay in the potential of the Milky Way is about $240$\,days; if path lengthening acted on light
only, the light would arrive about $120$\,days late \KM. Third, binary stars emit quadrupole radiation:
the orbital decay of the Hulse--Taylor pulsar follows the quadrupole formula,
$\dot P_{\text{GR}}=-2.402\times10^{-12}$, to $0.2\,\%$ \KM. The polarisation favours pure tensor modes
over pure scalar modes with a Bayes factor above $1000$ \QM. A scalar $m(x)$ can travel at $c$ but
supplies no tensor polarisation; the flat representation needs a tensor part for waves \SM. In the curvature representation it is
given with the adopted field equations (Doc.~312).

% ============================================================
\section{The speed-of-light limit and the field constants}
% ============================================================

In the basic form with $c=1$, $c$ is the conversion factor between the spatial directions and the
compact fourth direction $x_4$, not the speed of a medium (Doc.~134). The spectrum of the carrier is
$E^2=p^2+m^2$ with $m=k_4/R_4$; the winding number is the mass number (Doc.~314). The group velocity is
\begin{equation}
v=\frac{dE}{dp}=\frac{p}{\sqrt{p^2+m^2}}\quad\BM.
\end{equation}
It equals $1$ exactly when the mode has no winding in the time direction ($k_4=0$), and it is smaller
than $1$ for every massive mode at every momentum; the energy needed to reach $v=1$ diverges \BM. A
massive mode distributes its wave number between space and $x_4$, and only the spatial share is
motion. In the mass reading this means: whatever carries a clock is slower than light; only what has
no proper time reaches the limit (Doc.~391). The kinetic form of $\tilde T\cdot m=1$ is de Broglie's
harmony of phases: the inner frequency drops by $\gamma$, the phase frequency rises by $\gamma$, and the
product stays (Doc.~312) \BM.

The limit is the same in all representations. A conformal rescaling $\Omega^2\eta_{\mu\nu}$ leaves the
null vectors fixed, and the locally measured ratio of proper length to proper time of a light signal
does not depend on $\Omega$ \BM; mass variation can never change $c$ locally. What changes in a field
is the coordinate speed: it is $1$ for pure mass variation and $1-2U=1/n$ with $n=1+2U$ in the
isotropic field with $\gamma=1$, smaller by $4.2\times10^{-6}$ at the solar limb \KM. This slowing
belongs to the spatial half, the path lengthening, not to a variable $c$ (R132).

The same holds for the electric and the magnetic field constant. A static field acts on Maxwell fields
like a medium with $\varepsilon=\mu=\sqrt{B/A}$ for $g_{00}=-A$, $g_{ij}=B\,\delta_{ij}$
(Plebanski~\cite{Plebanski1960}) \QM. For $\gamma=1$ both rise equally by $1+2U$, the coordinate speed
$1/\sqrt{\varepsilon\mu}$ drops by $2U$, and the wave impedance $\sqrt{\mu/\varepsilon}=1$ stays \BM.
This is why light in a gravitational field is neither reflected nor deflected birefringently, and why
the deflection depends neither on colour nor on polarisation. Pure mass variation leaves
$\varepsilon=\mu=1$, because Maxwell's equations are conformally invariant in four dimensions \BM;
here too the spatial half carries the whole effect, in equal parts in $\varepsilon$ and $\mu$. Measured
locally, $\varepsilon_0$, $\mu_0$ and $\varepsilon_0\mu_0c^2=1$ remain the same everywhere; in the
Heaviside--Lorentz units of the basic form both are one. The frame dragging of Section~6 appears in
this medium form as a magnetoelectric coupling $\vec w=-g_{0i}/g_{00}$ between the electric and the
magnetic field (Plebanski) \QM.

% ============================================================
\section{Result}
% ============================================================

\begin{center}
\begin{tabular}{@{}p{4.2cm}p{5.4cm}p{3.6cm}@{}}
\toprule
Domain & Mass reading & Status \\
\midrule
Free fall, redshift, clocks & carries it alone, common core & \KM \\
Universe at large & exact (Weyl $=0$), passes three discriminators, AP as in FRW & \BM/\KM \\
distance law $1+z=e^{\xi x}$ & fails the AP test at $z=2.33$ ($5.3\,\sigma$) & \XM \\
Light, Mercury, geodetic & needs the spatial half, $\gamma=\beta=1$ & \KM, origin \SM \\
Rotating masses & needs a co-moving vector field & \BM, origin \SM \\
Gravitational waves & needs a tensor part & \SM \\
Speed of light, $\varepsilon_0$, $\mu_0$ & unchanged locally, light cones conformally invariant & \BM \\
\bottomrule
\end{tabular}
\end{center}

The three representations of Doc.~389 are equivalent for the conformally flat part and for the
static, non-rotating weak field with $\gamma=\beta=1$. The mass reading is the complete representation
of the conformally flat part: it carries the temporal half everywhere and the universe at large
exactly. What shows itself as space curvature in the narrower sense is the Weyl part around individual
masses; the flat representation has to supply it in addition, as the spatial half with the ruler sign
of proper lengths, as a co-moving field for rotation, and as a tensor part for waves. Mixed forms are
admissible as long as the spatial half carries the share one in total.

% ============================================================
\section{What remains open}
% ============================================================

\begin{itemize}[leftmargin=*,itemsep=1pt]
\item The origin of the spatial half from the structure of the carrier. Doc.~308 identifies the K3
area invariance as the only one of four tested candidates that yields $\gamma=1$ and confirms the
factor two of light deflection numerically; Doc.~308 itself lists the forward proof as open, as does
Doc.~389 \SM.
\item Dragging and the tensor part of gravitational waves in the flat representation. In the
curvature representation both are contained, because FFGFT adopts the field equations, whose form is
enforced by Lovelock's theorem (Doc.~312). Open is whether the structure of the carrier itself
supplies the co-moving field $\vec u$ and a tensor part. The corpus contains only scalar approaches
for this (Docs.~131, 189); the attempt in Doc.~143 to obtain the tensor part from the scalar is
circular according to Doc.~144 \SM.
\item A distance law of FFGFT that passes the Alcock--Paczyński test. The reciprocity
$d_A=d_L/(1+z)^2$ is no longer open: it holds exactly in reading M and is the same statement as $k=4$.
The pure exponential form $1+z=e^{\xi x}$ (Docs.~280, 308), read as coordinate distance of reading M,
fails at the Lyman-$\alpha$ value of DESI (Section~5). Open is which relation $z(d)$ the structure of
the carrier supplies and whether a different reading of the path length $x$ changes the result;
Doc.~268 so far treats $D_A$ as an external quantity \SM.
\end{itemize}

\vspace{1em}
\noindent\textit{Verification scripts:} \texttt{2/python/Dok392\_Skripte/} --- \texttt{pruef\_392a\_schwachfeld.py}
(14/14), \texttt{pruef\_392b\_mitfuehrung.py} (11/11), \texttt{pruef\_392c\_gravitationswellen.py} (8/8),
\texttt{pruef\_392d\_kosmos.py} (9/9), \texttt{pruef\_392e\_massenlesart.py} (11/11), \texttt{pruef\_392f\_lichtgrenze.py} (16/16), \texttt{pruef\_392g\_abstaende.py} (8/8),
\texttt{pruef\_392h\_desi\_ap.py} (4/4); together \texttt{run\_all\_392.sh}, 81/81.

\section*{Tool and method}
\addcontentsline{toc}{section}{Tool and method}

Claude (Anthropic) was used as a tool for: searching the corpus, writing and running the verification
scripts, preparing the \LaTeX\ sources according to the author's specifications. Every number in the
text is secured by a check in the scripts named; statements that only reproduce literature values are
marked \QM\ and are not counted there as checks.

\begin{thebibliography}{99}
\bibitem{Bertotti2003} B.~Bertotti, L.~Iess, P.~Tortora, \textit{A test of general relativity using radio links with the Cassini spacecraft}, Nature \textbf{425} (2003) 374--376.
\bibitem{Everitt2011} C.~W.~F.~Everitt et al., \textit{Gravity Probe B: Final results of a space experiment to test general relativity}, Phys.\ Rev.\ Lett.\ \textbf{106} (2011) 221101.
\bibitem{Williams2012} J.~G.~Williams, S.~G.~Turyshev, D.~H.~Boggs, \textit{Lunar laser ranging tests of the equivalence principle}, Class.\ Quantum Grav.\ \textbf{29} (2012) 184004.
\bibitem{Ciufolini2004} I.~Ciufolini, E.~C.~Pavlis, \textit{A confirmation of the general relativistic prediction of the Lense--Thirring effect}, Nature \textbf{431} (2004) 958--960.
\bibitem{Abbott2017a} B.~P.~Abbott et al., \textit{Gravitational waves and gamma-rays from a binary neutron star merger: GW170817 and GRB 170817A}, Astrophys.\ J.\ Lett.\ \textbf{848} (2017) L13.
\bibitem{Abbott2017b} B.~P.~Abbott et al., \textit{GW170814: A three-detector observation of gravitational waves from a binary black hole coalescence}, Phys.\ Rev.\ Lett.\ \textbf{119} (2017) 141101.
\bibitem{Weisberg2016} J.~M.~Weisberg, Y.~Huang, \textit{Relativistic measurements from timing the binary pulsar PSR B1913+16}, Astrophys.\ J.\ \textbf{829} (2016) 55.
\bibitem{Pound1960} R.~V.~Pound, G.~A.~Rebka, \textit{Apparent weight of photons}, Phys.\ Rev.\ Lett.\ \textbf{4} (1960) 337--341.
\bibitem{Chou2010} C.~W.~Chou et al., \textit{Optical clocks and relativity}, Science \textbf{329} (2010) 1630--1633.
\bibitem{Vessot1980} R.~F.~C.~Vessot et al., \textit{Test of relativistic gravitation with a space-borne hydrogen maser}, Phys.\ Rev.\ Lett.\ \textbf{45} (1980) 2081--2084.
\bibitem{Delva2018} P.~Delva et al., \textit{Gravitational redshift test using eccentric Galileo satellites}, Phys.\ Rev.\ Lett.\ \textbf{121} (2018) 231101.
\bibitem{DESI2025} M.~Abdul Karim et al.\ (DESI Collaboration), \textit{DESI DR2 Results II: Measurements of baryon acoustic oscillations and cosmological constraints}, arXiv:2503.14738 (2025), DOI 10.1103/tr6y-kpc6; Table~IV.
\bibitem{Wetterich2013} C.~Wetterich, \textit{A Universe without expansion}, Phys.\ Dark Univ.\ \textbf{2} (2013) 184--187.
\bibitem{White2024} R.~M.~T.~White et al., \textit{The Dark Energy Survey supernova program: slow supernovae show cosmological time dilation out to $z\sim1$}, Mon.\ Not.\ R.\ Astron.\ Soc.\ \textbf{533} (2024) 3365--3378.
\bibitem{Noterdaeme2011} P.~Noterdaeme et al., \textit{The evolution of the cosmic microwave background temperature}, Astron.\ Astrophys.\ \textbf{526} (2011) L7.
\bibitem{Lubin2001} L.~M.~Lubin, A.~Sandage, \textit{The Tolman surface brightness test for the reality of the expansion. IV.}, Astron.\ J.\ \textbf{122} (2001) 1084--1103.
\bibitem{Plebanski1960} J.~Plebanski, \textit{Electromagnetic waves in gravitational fields}, Phys.\ Rev.\ \textbf{118} (1960) 1396--1408.
\end{thebibliography}
